Between 2015 and 2024, CO$_2$ emissions from Chilean copper production fell by roughly 17\% despite a sustained increase in energy consumption. The decomposition analysis shows that this was not the result of technological improvements in mining itself. The carbon intensity of combustion fuels stayed roughly constant over the period, and the share of electricity in total energy use barely moved. What changed was the electricity grid: as coal was progressively replaced by solar and wind, the carbon intensity of electricity fell by over 37\%, pulling down indirect emissions with it.

This matters for how we interpret the current situation. The progress made so far has come largely for free from a mining perspective, driven by Chile's broader energy transition rather than by anything the sector has done differently. Direct emissions, which come mainly from diesel-powered machinery in open-pit operations, actually increased over the period. As ore grades continue to decline and extraction goes deeper, this pressure will only grow.

The path forward depends on two things happening together. First, the grid needs to keep decarbonizing. The IEA projects a carbon intensity of 29~gCO$_2$/kWh by 2035, down from around 189 today. At that level, electrifying haul trucks or concentrating plants would produce emission reductions, not just a change in where the emissions are counted. Second, the sector itself needs to start shifting away from diesel. That transition is expensive and slow replacing a mining fleet takes years and significant capital but section \ref{sec:opportunities} shows there are already viable technologies on the market. The two conditions reinforce each other: a cleaner grid makes the investment in electrification more worthwhile, and electrification makes the grid decarbonization more impactful for copper emissions.

Chile's geography is an asset in this respect. The northern desert, where most of the country's copper is mined, also has some of the highest solar irradiance in the world. The conditions for a low-carbon copper industry are better here than almost anywhere else. The question is whether the investment follows.